\documentclass[11pt]{amsart}
\usepackage[margin=1.15in]{geometry}
\usepackage{amsmath,amssymb,amsthm}
\usepackage[hidelinks]{hyperref}
\newtheorem{theorem}{Theorem}
\newtheorem{corollary}[theorem]{Corollary}
\newtheorem{lemma}[theorem]{Lemma}
\newtheorem{proposition}[theorem]{Proposition}
\theoremstyle{remark}
\newtheorem*{remark}{Remark}
\numberwithin{equation}{section}
\newcommand{\ZZ}{\mathbb Z}
\newcommand{\Dp}{\mathbb D^{+}}
\newcommand{\fl}[1]{\lfloor #1\rfloor}
\newcommand{\Sb}{\overline{\mathcal S}}
\title[Bernoulli polynomials with integral derivatives]{Kellner's conjectures on Bernoulli polynomials\\ with integral derivatives}
\author{Carlos Kuban}
\date{Draft, November 2026}
\begin{document}
\begin{abstract}
We prove that the derivative $B_n'(x)$ of the $n$th Bernoulli polynomial has integral coefficients if and only if
$n\in\{1,2,4,6,10,12,28,30,36,60\}$, as conjectured by Kellner. This and several related conjectures of Kellner and of Alekseyev follow from
another conjecture of Kellner, which we also prove: for every $n>192$ there is a prime $p>\sqrt n$ such that the sum of the base-$p$ digits of $n$ is at least~$p$.
The proof is short. A telescoping argument shows that almost all integers $m$ slightly larger than $\sqrt n$ have base-$m$ digit sum at least $m$;
this is combined with a classical explicit estimate for the number of primes in short intervals and a small computation.
\end{abstract}
\maketitle

\section{Results}

For integers $m\ge2$ and $n\ge1$ let $s_m(n)$ be the sum of the base-$m$ digits of $n$, and let
\[
\Dp_n=\prod_{\substack{p>\sqrt n\\ s_p(n)\ge p}}p .
\]
Here and below $p$ denotes a prime, and $\omega(m)$ is the number of primes dividing $m$. The product is finite, as $s_p(n)=n$ for $p>n$.
Kellner \cite[Conjectures 1 and 2]{K17} conjectured that $\Dp_n>1$ for all $n>192$ and that $\omega(\Dp_n)\sim\kappa\sqrt n/\log n$ for some constant $\kappa>1$.
The asymptotic formula was proved, with $\kappa=2$, by Bordell\`es, Luca, Moree and Shparlinski \cite{BLMS}, but their method does not lead to an explicit
bound beyond which $\Dp_n>1$; see \cite[p.~2]{K24}. We prove the first conjecture.

\begin{theorem}\label{thm:main}
For every $n>192$ there is a prime $p>\sqrt n$ with $s_p(n)\ge p$. More generally $\omega(\Dp_n)\ge k$ for all $n>b_k$, where
\[
\begin{array}{c|cccccccc}
k&1&2&3&4&5&6&7&8\\ \hline
b_k&192&1080&2014&2430&3360&5076&5796&11824
\end{array}
\]
and these values of $b_k$ are best possible.
\end{theorem}

The following statements are consequences. That the sets in Corollaries \ref{cor:first}--\ref{cor:higher} are finite was shown by Kellner \cite{K24}.

\begin{corollary}[{\cite[Conjecture 1]{K24}; OEIS A094960}]\label{cor:first}
We have $B_n'(x)\in\ZZ[x]$ if and only if
\[
n\in\{1,2,4,6,10,12,28,30,36,60\}.
\]
\end{corollary}

\begin{corollary}[{\cite[Conjecture 11]{K24}; OEIS A195441}]\label{cor:R}
The denominator of $B_n(x)-B_n$ equals the radical of $n+1$ if and only if $n\in\{3,5,8,9,11,27,29,35,59\}$.
\end{corollary}

For $k\ge1$ let $\Sb_k$ be the set of all $n\ge1$ for which the $k$th derivative $B_n^{(k)}(x)$ is in $\ZZ[x]$.

\begin{corollary}[{\cite[Conjecture 15]{K24}; OEIS A366169, A366186--A366188}]\label{cor:higher}
\begin{align*}
\Sb_2&=\{1\text{--}7,9\text{--}13,15,16,21,25,28\text{--}31,36,37,55,57,60,61,70,121,190\},\\
\Sb_3&=\{1\text{--}18,20\text{--}22,25,26,28\text{--}32,35\text{--}38,42,50,52,55\text{--}58,60\text{--}62,\\
&\qquad66,70\text{--}72,78,80,92,110,121,122,156,176,177,190,191,210,392\},
\end{align*}
where $a\text{--}b$ stands for all integers from $a$ to $b$. Moreover
\[
\begin{array}{c|cccccccccccc}
k&1&2&3&4&5&6&7&8&9&10&11&12\\ \hline
|\Sb_k|&10&29&59&88&125&164&211&258&310&368&433&502\\
\max\Sb_k&60&190&392&393&904&905&906&907&908&1089&2080&2387
\end{array}
\]
In particular the lists in OEIS A366187 $(k=4)$ and A366188 $(k=5)$ are complete.
\end{corollary}

\begin{corollary}[{\cite[Conjecture 1]{K17}}]\label{cor:K17}
Let $P_n$ be the largest prime $p$ with $s_p(n)\ge p$. For $n>192$ it exists, and $P_n>\sqrt n$ and $(n+1)/P_n<\sqrt n$.
\end{corollary}

\begin{corollary}[Conjecture of M.~Alekseyev, OEIS A094960]\label{cor:alek}
For every integer $k>78$ there is a prime $p\ge3$ such that the sum of the last two base-$p$ digits of $k$ is at least $p$.
\end{corollary}

\section{The counting lemma}

Let $n\ge1$, $x=\sqrt n$ and $s=\fl x$. For an integer $m>s$ put $q(m)=\fl{n/m}$. As $m^2\ge(s+1)^2>n$, we have $q(m)<m$, so $n$ has the base-$m$ digits $q(m)$ and $n-q(m)m$ and
\begin{equation}\label{eq:digits}
s_m(n)=n-(m-1)\,q(m).
\end{equation}
Let us call an integer $m>s$ \emph{good} if $s_m(n)\ge m$, and \emph{bad} otherwise.

\begin{lemma}\label{lem:one}
Let $m\ge s+2$. Then $q(m-1)-q(m)\in\{0,1\}$, and $m$ is bad if and only if $q(m-1)=q(m)$ or $m-1$ divides $n$.
\end{lemma}
\begin{proof}
We have $0\le q(m-1)-q(m)<\frac n{m-1}-\frac nm+1=\frac n{m(m-1)}+1<2$, because $m(m-1)>(s+1)^2>n$.
By \eqref{eq:digits}, $m$ is good if and only if $n-1\ge(m-1)(q(m)+1)$, that is, $\fl{\frac{n-1}{m-1}}>q(m)$ (see \cite[Lemma 6]{K17}). Now
$\fl{\frac{n-1}{m-1}}$ is $q(m-1)-1$ if $m-1$ divides $n$, and $q(m-1)$ otherwise. As $q(m-1)-q(m)\le1$, the claim follows.
\end{proof}

\begin{lemma}\label{lem:bad}
For every $U\ge1$ at most $4+\frac{2U^2}{s+U}$ of the integers $s+1,\dots,s+U$ are bad.
\end{lemma}
\begin{proof}
Put $v=s+U-x$; then $0<v\le U$ and $x+v=s+U$. We use
\begin{equation}\label{eq:frac}
\frac n{s+U}=\frac{x^2}{x+v}=x-v+\frac{v^2}{x+v}\le x-v+\frac{U^2}{s+U}.
\end{equation}

(a) By Lemma \ref{lem:one} each of the differences $q(m-1)-q(m)$, $s+2\le m\le s+U$, is $0$ or $1$, and their sum is $q(s+1)-q(s+U)$. Hence the number of these $m$ with
$q(m-1)=q(m)$ is $U-1-q(s+1)+q(s+U)$. We have $q(s+1)\ge s-1$, as $(s-1)(s+1)<n$, and \eqref{eq:frac} bounds $q(s+U)$. So the number in question is at most
\[
U-s+x-v+\frac{U^2}{s+U}=2(x-s)+\frac{U^2}{s+U}<2+\frac{U^2}{s+U}.
\]

(b) Let $d_1<\dots<d_h$ be the divisors of $n$ in $[s+1,s+U-1]$. The function $y+n/y$ is at least $2x$ and strictly increasing for $y\ge x$. So the numbers $t_i=d_i+n/d_i$ are integers with
$2x\le t_1<\dots<t_h\le s+U+\frac n{s+U}\le2x+\frac{U^2}{s+U}$, by \eqref{eq:frac}. Hence $h\le1+\frac{U^2}{s+U}$.

By Lemma \ref{lem:one}, a bad $m$ with $s+2\le m\le s+U$ is counted in (a), or $m-1$ is one of the divisors in (b). Together with $m=s+1$ this gives the bound.
\end{proof}

The number of good integers was bounded from above, by $\sqrt n$, in \cite[Lemma 7]{K17}; here we need that there are many of them.
If a prime $p>s$ is good, then $p\mid\Dp_n$. So for all $n\ge1$ and $U\ge1$
\begin{equation}\label{eq:main}
\omega(\Dp_n)\ge\pi(s+U)-\pi(s)-4-\frac{2U^2}{s+U}.
\end{equation}

\section{Large \texorpdfstring{$n$}{n}}

Let $s\ge10^9$ and write $\ell=\log s>20$. With $\vartheta(y)=\sum_{p\le y}\log p$, Rosser and Schoenfeld \cite{RS} (see \cite[Table 2]{BKLNW}) proved
\begin{equation}\label{eq:RS}
|\vartheta(y)-y|<\frac y{40\log y}\qquad(y\ge678\,407).
\end{equation}
Much stronger bounds are known; for instance $|\vartheta(y)-y|<0.2\,y/\log^2y$ for $y\ge3\,594\,641$ \cite[Theorem 4.2]{D18}. Take $U=\fl{s/(4\ell)}$. Then $U\le s/80$ and $\log(s+U)\le\ell+\frac1{4\ell}\le1.00063\,\ell$, so
\[
\vartheta(s+U)-\vartheta(s)\ge U-\frac{2s+U}{40\,\ell}\ge\frac s{4\ell}-1-0.0504\,\frac s\ell=0.1996\,\frac s\ell-1,
\]
\[
\pi(s+U)-\pi(s)\ge\frac{\vartheta(s+U)-\vartheta(s)}{\log(s+U)}\ge0.1994\,\frac s{\ell^2}-1 .
\]
Since $\frac{2U^2}{s+U}\le\frac s{8\ell^2}$, \eqref{eq:main} gives
\begin{equation}\label{eq:large}
\omega(\Dp_n)\ge0.0744\,\frac s{\log^2s}-5>10^5\qquad(s=\fl{\sqrt n}\ge10^9),
\end{equation}
because $s/\log^2s$ is increasing for $s>e^2$ and exceeds $2.3\cdot10^6$ at $s=10^9$.

\section{Small \texorpdfstring{$n$}{n}}

It remains to consider $n<10^{18}$. A prime $p>s$ is bad for a given $n\in[s^2,s^2+2s]$ if and only if, with $q=\fl{n/p}$, we have $q+(n-qp)<p$. So the $n$ for which $p$ is bad are those in the intervals
\[
[\,qp,\ qp+p-1-q\,]\cap[s^2,s^2+2s]\qquad(q\ge0),
\]
and only a few of these are non-empty. Given a list of primes exceeding $s$, the largest number of them that are bad for the same $n\in[s^2,s^2+2s]$ is therefore found by sorting the end points of these intervals.
We did this for each $s$ with $1500\le s<10^9$ and the first $40$ primes exceeding $s$. It turned out that never more than $16$ of them are bad for the same $n$; hence
\begin{equation}\label{eq:middle}
\omega(\Dp_n)\ge24\qquad(1500^2\le n<10^{18}).
\end{equation}
For $n<1500^2$ we computed $\omega(\Dp_n)$ from the definition, as far as needed; the largest $n$ with $\omega(\Dp_n)<k$ is $b_k$, for $k\le8$. With \eqref{eq:large} and \eqref{eq:middle} this proves Theorem \ref{thm:main}.

The computation uses integer arithmetic only. It takes about $25$ minutes on one processor core, almost all of it for $s\ge4\cdot10^6$, a range which could also be covered by \eqref{eq:RS}. For $1500\le s\le4\cdot10^6$ it was repeated with a second, independently written program.
The $n$ with $\Dp_n=1$ are $1$, $2$, $4$, $6$, $10$, $11$, $12$, $15$, $16$, $28$, $29$, $30$, $35$, $36$, $58$, $59$, $60$, $69$, $70$, $78$, $79$, $80$, $174$, $190$, $191$, $192$.

\section{Denominators}

The following facts are due to Kellner and Sondow \cite{KS} and Kellner \cite{K17,K24}. We include a short proof. Let $(n)_k=n(n-1)\cdots(n-k+1)$.

\begin{proposition}\label{prop:den}
Let $n\ge1$. The denominators of $B_n(x)-B_n$ and of $B_n^{(k)}(x)$, $1\le k\le n$, are squarefree. A prime $p$ divides the first of them if and only if $s_p(n)\ge p$, and the second if and only if
$p\nmid(n)_k$ and $s_p(n-k+1)\ge p$.
\end{proposition}
\begin{proof}
By the theorem of von Staudt and Clausen $B_j$ is $p$-integral for $j\notin J_p$ and has $p$-adic valuation $-1$ for $j\in J_p$, where $J_p$ is the set of the positive multiples of $p-1$ if $p\ge3$, and $J_2=\{1,2,4,6,8,\dots\}$.
Write $j\preceq N$ if no base-$p$ digit of $j$ exceeds the corresponding digit of $N$. By Lucas's theorem $p\nmid\binom Nj$ if and only if $j\preceq N$. In this case $s_p(j)\le s_p(N)$, with equality only for $j=N$. We claim that, for $N\ge0$,
\begin{itemize}
\item[(a)] there is $j\in J_p$ with $j\preceq N$ if and only if $s_p(N)\ge p-1$;
\item[(b)] there is $j\in J_p$ with $j\preceq N$ and $j\ne N$ if and only if $s_p(N)\ge p$.
\end{itemize}
Indeed, if $j\in J_p$, then $s_p(j)\equiv j\equiv0\pmod{p-1}$ and $s_p(j)>0$, so $s_p(j)\ge p-1$; this shows that the conditions are necessary. If $s_p(N)\ge p-1$ we can choose $j\preceq N$ with $s_p(j)=p-1$. Then $j$ is a positive multiple of $p-1$, and for $p=2$ it is a power of $2$. So $j\in J_p$, and if $s_p(N)\ge p$ then $j\ne N$.

Now $B_n(x)-B_n=\sum_{j<n}\binom njB_jx^{n-j}$, and the first claim follows from (b) with $N=n$. Also
\[
B_n^{(k)}(x)=(n)_kB_{n-k}(x)=\sum_{j\le n-k}(n)_k\binom{n-k}jB_jx^{n-k-j}.
\]
If $p\mid(n)_k$, its coefficients are $p$-integral. If $p\nmid(n)_k$, then by (a) they are not all $p$-integral if and only if $s_p(n-k)\ge p-1$, and $p\nmid n-k+1$ implies $s_p(n-k+1)=s_p(n-k)+1$.
\end{proof}

\section{Proofs of the corollaries}

Let $m\ge1$ and $s=\fl{\sqrt m}$. We use two remarks.
\begin{itemize}
\item[(A)] If $p\mid\Dp_m$ then $p\nmid m$. Indeed, if $p>\sqrt m$ and $p\mid m$, then $s_p(m)=m/p<p$.
\item[(B)] If $0\le j\le s+1$ then at most one prime $p>\sqrt m$ divides $m+j$, since two such primes have product at least $(s+1)(s+2)>s^2+2s+j\ge m+j$.
\end{itemize}

\begin{proof}[Proof of Corollary \ref{cor:first}]
By Proposition \ref{prop:den}, $B_n'(x)\in\ZZ[x]$ if and only if there is no prime $p\nmid n$ with $s_p(n)\ge p$. By Theorem \ref{thm:main} and (A) there is one if $n>192$. The $n\le192$ were checked directly.
\end{proof}

\begin{proof}[Proof of Corollary \ref{cor:R}]
If the denominator of $B_n(x)-B_n$ equals the radical of $n+1$, then by Proposition \ref{prop:den} all primes dividing $\Dp_n$ divide $n+1$. So $\omega(\Dp_n)\le1$ by (B), and $n\le1080$ by Theorem \ref{thm:main}. These $n$ were checked directly.
\end{proof}

\begin{proof}[Proof of Corollary \ref{cor:higher}]
All $n<k$ belong to $\Sb_k$. For $n\ge k$ put $m=n-k+1$. By Proposition \ref{prop:den}, $n\in\Sb_k$ if and only if every prime $p$ with $s_p(m)\ge p$ divides $m(m+1)\cdots(m+k-1)$.
Suppose that $k\le\fl{\sqrt m}+2$ and $\omega(\Dp_m)\ge k$. By (A) no prime dividing $\Dp_m$ divides $m$, and by (B) each of $m+1,\dots,m+k-1$ is divisible by at most one of them, so that $n\notin\Sb_k$.
By \eqref{eq:large} and \eqref{eq:middle} this applies, for all $k\le24$, to all $m\ge1500^2$ (and for $k\le8$, by Theorem \ref{thm:main}, to all $m>b_k$). The remaining $m$ were checked.
\end{proof}

\begin{proof}[Proof of Corollary \ref{cor:K17}]
The first claim is Theorem \ref{thm:main}. If $n>1080$ then $\omega(\Dp_n)\ge2$, hence $P_n\ge\fl{\sqrt n}+2>\sqrt n+1/\sqrt n=(n+1)/\sqrt n$. For $192<n\le1080$ the second claim was checked directly.
\end{proof}

\begin{proof}[Proof of Corollary \ref{cor:alek}]
If $k>192$ take $p>\sqrt k$ as in Theorem \ref{thm:main}; then $k$ has at most two digits in base $p$, and $p\ge3$. For $78<k\le192$ the claim was checked directly.
\end{proof}

\begin{remark}
The small cases were checked with Proposition \ref{prop:den}. For Corollary \ref{cor:R}, and for Corollaries \ref{cor:first} and \ref{cor:higher} with $k\le6$, they were also confirmed from exact values of the Bernoulli numbers, without using Proposition \ref{prop:den}; by Theorem \ref{thm:main} and the proofs above it suffices to consider $n\le1080$ and $n\le5081$, respectively.
\end{remark}

\section{Formal verification}

All arguments of this paper except the estimate \eqref{eq:RS} have been formalized in Lean~4 with Mathlib (version~4.36.0); the code accompanies the paper.
The estimate \eqref{eq:RS} is not available in Mathlib. It enters the formal proof as an explicit hypothesis, required only for $y\ge4\cdot10^6$, and not as an axiom; the formal statements are of the form ``if \eqref{eq:RS} holds for $y\ge4\cdot10^6$, then \dots''. Apart from this hypothesis the formal proofs use only the standard axioms of Lean.
In the formalization the constant $4$ in Lemma~\ref{lem:bad} is replaced by $6$, Proposition~\ref{prop:den} is derived from the theorem of von Staudt and Clausen and from Kummer's theorem as they are stated in Mathlib, and every computation is carried out by Lean's kernel (no compiled code is trusted). For $n<1.6\cdot10^{13}$ this computation shows that there are at least three primes $p>\sqrt n$ with $s_p(n)\ge p$ whenever $n\ge2015$; so Theorem~\ref{thm:main} is formally verified without any hypothesis for $n<1.6\cdot10^{13}$.
The formal version of Corollary~\ref{cor:R} is stated exactly as the statement \texttt{OeisA195441.conjecture} of the \texttt{formal-conjectures} project.

\subsection*{Acknowledgement}
This work was carried out with the assistance of a large language model (Claude, by Anthropic), which was used for the proof, the computations, the Lean formalization and the writing. The paper has not yet been refereed. The programs and the Lean code are available with it.

\begin{thebibliography}{9}
\bibitem{BLMS} O. Bordell\`es, F. Luca, P. Moree, I. E. Shparlinski, \emph{Denominators of Bernoulli polynomials}, Mathematika 64 (2018), 519--541.
\bibitem{BKLNW} S. Broadbent, H. Kadiri, A. Lumley, N. Ng, K. Wilk, \emph{Sharper bounds for the Chebyshev function $\theta(x)$}, Math. Comp. 90 (2021), 2281--2315; arXiv:2002.11068.
\bibitem{D18} P. Dusart, \emph{Explicit estimates of some functions over primes}, Ramanujan J. 45 (2018), 227--251.
\bibitem{K17} B. C. Kellner, \emph{On a product of certain primes}, J. Number Theory 179 (2017), 126--141; arXiv:1705.04303.
\bibitem{K24} B. C. Kellner, \emph{On the finiteness of Bernoulli polynomials whose derivative has only integral coefficients}, J. Integer Seq. 27 (2024), Article 24.2.8; arXiv:2310.01325.
\bibitem{KS} B. C. Kellner, J. Sondow, \emph{Power-sum denominators}, Amer. Math. Monthly 124 (2017), 695--709; arXiv:1705.03857.
\bibitem{RS} J. B. Rosser, L. Schoenfeld, \emph{Sharper bounds for the Chebyshev functions $\theta(x)$ and $\psi(x)$}, Math. Comp. 29 (1975), 243--269.
\end{thebibliography}
\end{document}
